\section{系统局限性与后续技术演进路线}
\label{sec:limitations}

\subsection{交互总量高流量截断偏差及其优化方向}
在当前基准测试中，交互总量（Interaction Volume）的全局截断阈值固定设置为 $C=20$。由于真实企业环境存在大量依赖自动化工具或深度集成 IDE 的高频研发人员，固定截断导致整体交互总量呈现约 $39.8\%$ 的结构性低估（Under-counting Bias）。

后续技术攻关的核心在于研究\textbf{截断阈值 $C$ 与隐私参数 $\varepsilon$ 的联合优化模型（$C \times \varepsilon$ Joint Calibration）}：
\begin{equation}
  \min_{C} \mathbb{E}\left[ \left( \hat{M}(C) - M \right)^2 \right] = \min_{C} \left[ \text{Bias}^2(C) + \text{Var}(C, \varepsilon) \right]
\end{equation}
即在“截断造成的负向偏差”与“增大 $C$ 导致敏感度上升从而放大的高斯方差”之间寻找理论最优平衡点。

\subsection{自适应贡献截断的隐私理论考量}
针对截断阈值 $C$ 的选择，一种工程直觉是采用“基于数据分布自适应选择分位数（如 90\% 分位数）”。然而，在严格的差分隐私理论下必须清醒认识到：\textbf{直接从待保护的数据集自身计算分位数并将阈值反馈用于后续截断，该过程本身构成了对私有数据的观测，若不分配隐私预算将产生严重的隐私泄露}。

因此，后续关于自适应截断（Adaptive Clipping）的演进路线规划如下：
\begin{enumerate}[label=(\alph*), leftmargin=2.5em]
  \item \textbf{独立先验校准集方案：} 利用历史已过期、不属于当前敏感窗口的公共或历史交互分布估算相对稳定的经验阈值 $C$；
  \item \textbf{差分隐私私有分位数分配方案：} 将单次查询的总预算拆分为两部分：分配 $\varepsilon_{\text{clip}}$ 运行满足差分隐私的私有分位数估计机制，求出加噪截断阈值 $\tilde{C}$；随后分配 $\varepsilon_{\text{sum}} = \varepsilon - \varepsilon_{\text{clip}}$ 执行截断后的高斯聚合。该方案仍需在效用收益与分段预算损耗之间开展系统性实证评估。
\end{enumerate}

\subsection{细粒度业务分类体系的拓展适配}
当前系统的基准分类体系包含 3 个核心任务（Technical、Document、Analysis）。在企业实际生产场景中，业务部门往往期望进一步拆分为数十个子类别（例如代码审查、单元测试生成、SQL 编写、长文润色、合规审查等）。

随着分类维度 $D$ 的扩张，稀疏单元（Sparsity）数量显著增加。在微小调用单元中，高斯机制的随机扰动相对占比将增大。后续研发将探索树状层级分类法（Hierarchical Taxonomy）与分层差分隐私投影算法，在宏观层级保持高精度的同时，在微观叶子层级实施自适应样本抑制。

\subsection{更多业务维度的扩展规划}
在当前五类核心统计指标的基础上，后续版本将逐步演进支持更多维度的指标 Profile：
\begin{itemize}[leftmargin=2em]
  \item \textbf{Token 吞吐量统计（Token Volume Profile）：} 针对单次交互中数千乃至数万的输入/输出 Token 量，推导基于重尾对数截断的差分隐私释放方案；
  \item \textbf{推理延迟与 SLA 达标率（Latency Profile）：} 评估不同供应商在峰值负载下的服务响应时间分布；
  \item \textbf{错误与降级重试率（Fallback Rate Profile）：} 评估模型网关主备切换与重试频次，指导高可用架构调优。
\end{itemize}

\subsection{智能化效用规划器 (Utility-Aware Planner) 的演进}
当前版本已实现了基础的离线效用对比、误差拆解与 $\varepsilon$-Sweep 模拟。下一阶段的研发目标是将规划器（Planner）从“离线辅助分析”演进为“智能化闭环推荐”：即系统根据输入数据集规模 $N$、历史活跃特征以及目标接收方的剩余隐私配额，自动为用户推荐兼顾合规底线与数据可用性的候选发布时间窗口、推荐 $\varepsilon$ 建议值以及配套的截断参数建议。
